\documentclass[a4paper,12pt,twoside]{memoir}

% Castellano
\usepackage[spanish,es-tabla]{babel}
\selectlanguage{spanish}
\usepackage[utf8]{inputenc}
\usepackage[T1]{fontenc}
\usepackage{lmodern} % Scalable font
\usepackage{microtype}
\usepackage{placeins}

\RequirePackage{booktabs}
\RequirePackage[table]{xcolor}
\RequirePackage{xtab}
\RequirePackage{multirow}

% Links
\usepackage[colorlinks]{hyperref}
\hypersetup{
	allcolors = {red}
}

% Ecuaciones
\usepackage{amsmath}

% Rutas de fichero / paquete
\newcommand{\ruta}[1]{{\sffamily #1}}

% Párrafos
\nonzeroparskip


% Imagenes
\usepackage{graphicx}
\newcommand{\imagen}[2]{
	\begin{figure}[!h]
		\centering
		\includegraphics[width=0.9\textwidth]{#1}
		\caption{#2}\label{fig:#1}
	\end{figure}
	\FloatBarrier
}

\newcommand{\imagenflotante}[2]{
	\begin{figure}%[!h]
		\centering
		\includegraphics[width=0.9\textwidth]{#1}
		\caption{#2}\label{fig:#1}
	\end{figure}
}



% El comando \figura nos permite insertar figuras comodamente, y utilizando
% siempre el mismo formato. Los parametros son:
% 1 -> Porcentaje del ancho de página que ocupará la figura (de 0 a 1)
% 2 --> Fichero de la imagen
% 3 --> Texto a pie de imagen
% 4 --> Etiqueta (label) para referencias
% 5 --> Opciones que queramos pasarle al \includegraphics
% 6 --> Opciones de posicionamiento a pasarle a \begin{figure}
\newcommand{\figuraConPosicion}[6]{%
  \setlength{\anchoFloat}{#1\textwidth}%
  \addtolength{\anchoFloat}{-4\fboxsep}%
  \setlength{\anchoFigura}{\anchoFloat}%
  \begin{figure}[#6]
    \begin{center}%
      \Ovalbox{%
        \begin{minipage}{\anchoFloat}%
          \begin{center}%
            \includegraphics[width=\anchoFigura,#5]{#2}%
            \caption{#3}%
            \label{#4}%
          \end{center}%
        \end{minipage}
      }%
    \end{center}%
  \end{figure}%
}

%
% Comando para incluir imágenes en formato apaisado (sin marco).
\newcommand{\figuraApaisadaSinMarco}[5]{%
  \begin{figure}%
    \begin{center}%
    \includegraphics[angle=90,height=#1\textheight,#5]{#2}%
    \caption{#3}%
    \label{#4}%
    \end{center}%
  \end{figure}%
}
% Para las tablas
\newcommand{\otoprule}{\midrule [\heavyrulewidth]}
%
% Nuevo comando para tablas pequeñas (menos de una página).
\newcommand{\tablaSmall}[5]{%
 \begin{table}
  \begin{center}
   \rowcolors {2}{gray!35}{}
   \begin{tabular}{#2}
    \toprule
    #4
    \otoprule
    #5
    \bottomrule
   \end{tabular}
   \caption{#1}
   \label{tabla:#3}
  \end{center}
 \end{table}
}

%
% Nuevo comando para tablas pequeñas (menos de una página).
\newcommand{\tablaSmallSinColores}[5]{%
 \begin{table}[H]
  \begin{center}
   \begin{tabular}{#2}
    \toprule
    #4
    \otoprule
    #5
    \bottomrule
   \end{tabular}
   \caption{#1}
   \label{tabla:#3}
  \end{center}
 \end{table}
}

\newcommand{\tablaApaisadaSmall}[5]{%
\begin{landscape}
  \begin{table}
   \begin{center}
    \rowcolors {2}{gray!35}{}
    \begin{tabular}{#2}
     \toprule
     #4
     \otoprule
     #5
     \bottomrule
    \end{tabular}
    \caption{#1}
    \label{tabla:#3}
   \end{center}
  \end{table}
\end{landscape}
}

%
% Nuevo comando para tablas grandes con cabecera y filas alternas coloreadas en gris.
\newcommand{\tabla}[6]{%
  \begin{center}
    \tablefirsthead{
      \toprule
      #5
      \otoprule
    }
    \tablehead{
      \multicolumn{#3}{l}{\small\sl continúa desde la página anterior}\\
      \toprule
      #5
      \otoprule
    }
    \tabletail{
      \hline
      \multicolumn{#3}{r}{\small\sl continúa en la página siguiente}\\
    }
    \tablelasttail{
      \hline
    }
    \bottomcaption{#1}
    \rowcolors {2}{gray!35}{}
    \begin{xtabular}{#2}
      #6
      \bottomrule
    \end{xtabular}
    \label{tabla:#4}
  \end{center}
}

%
% Nuevo comando para tablas grandes con cabecera.
\newcommand{\tablaSinColores}[6]{%
  \begin{center}
    \tablefirsthead{
      \toprule
      #5
      \otoprule
    }
    \tablehead{
      \multicolumn{#3}{l}{\small\sl continúa desde la página anterior}\\
      \toprule
      #5
      \otoprule
    }
    \tabletail{
      \hline
      \multicolumn{#3}{r}{\small\sl continúa en la página siguiente}\\
    }
    \tablelasttail{
      \hline
    }
    \bottomcaption{#1}
    \begin{xtabular}{#2}
      #6
      \bottomrule
    \end{xtabular}
    \label{tabla:#4}
  \end{center}
}

%
% Nuevo comando para tablas grandes sin cabecera.
\newcommand{\tablaSinCabecera}[5]{%
  \begin{center}
    \tablefirsthead{
      \toprule
    }
    \tablehead{
      \multicolumn{#3}{l}{\small\sl continúa desde la página anterior}\\
      \hline
    }
    \tabletail{
      \hline
      \multicolumn{#3}{r}{\small\sl continúa en la página siguiente}\\
    }
    \tablelasttail{
      \hline
    }
    \bottomcaption{#1}
  \begin{xtabular}{#2}
    #5
   \bottomrule
  \end{xtabular}
  \label{tabla:#4}
  \end{center}
}



\definecolor{cgoLight}{HTML}{EEEEEE}
\definecolor{cgoExtralight}{HTML}{FFFFFF}

%
% Nuevo comando para tablas grandes sin cabecera.
\newcommand{\tablaSinCabeceraConBandas}[5]{%
  \begin{center}
    \tablefirsthead{
      \toprule
    }
    \tablehead{
      \multicolumn{#3}{l}{\small\sl continúa desde la página anterior}\\
      \hline
    }
    \tabletail{
      \hline
      \multicolumn{#3}{r}{\small\sl continúa en la página siguiente}\\
    }
    \tablelasttail{
      \hline
    }
    \bottomcaption{#1}
    \rowcolors[]{1}{cgoExtralight}{cgoLight}

  \begin{xtabular}{#2}
    #5
   \bottomrule
  \end{xtabular}
  \label{tabla:#4}
  \end{center}
}


















\graphicspath{ {./img/} }

% Capítulos
\chapterstyle{bianchi}
\newcommand{\capitulo}[2]{
	\setcounter{chapter}{#1}
	\setcounter{section}{0}
	\chapter*{#2}
	\addcontentsline{toc}{chapter}{#2}
	\markboth{#2}{#2}
}

% Apéndices
\renewcommand{\appendixname}{Apéndice}
\renewcommand*\cftappendixname{\appendixname}

\newcommand{\apendice}[1]{
	%\renewcommand{\thechapter}{A}
	\chapter{#1}
}

\renewcommand*\cftappendixname{\appendixname\ }

% Formato de portada
\makeatletter
\usepackage{xcolor}
\newcommand{\tutor}[1]{\def\@tutor{#1}}
\newcommand{\course}[1]{\def\@course{#1}}
\definecolor{cpardoBox}{HTML}{E6E6FF}
\def\maketitle{
  \null
  \thispagestyle{empty}
  % Cabecera ----------------
\noindent\includegraphics[width=\textwidth]{cabecera}\vspace{1cm}%
  \vfill
  % Título proyecto y escudo informática ----------------
  \colorbox{cpardoBox}{%
    \begin{minipage}{.8\textwidth}
      \vspace{.5cm}\Large
      \begin{center}
      \textbf{TFG del Grado en Ingeniería Informática}\vspace{.6cm}\\
      \textbf{\LARGE\@title{}}
      \end{center}
      \vspace{.2cm}
    \end{minipage}

  }%
  \hfill\begin{minipage}{.20\textwidth}
    \includegraphics[width=\textwidth]{escudoInfor}
  \end{minipage}
  \vfill
  % Datos de alumno, curso y tutores ------------------
  \begin{center}%
  {%
    \noindent\LARGE
    Presentado por \@author{}\\ 
    en Universidad de Burgos --- \@date{}\\
    Tutor: \@tutor{}\\
  }%
  \end{center}%
  \null
  \cleardoublepage
  }
\makeatother

\newcommand{\nombre}{Bárbara Sainz Crespo} %%% cambio de comando

% Datos de portada
\title{título del TFG}
\author{\nombre}
\tutor{nombre tutor}
\date{\today}

\begin{document}

\maketitle


\newpage\null\thispagestyle{empty}\newpage


%%%%%%%%%%%%%%%%%%%%%%%%%%%%%%%%%%%%%%%%%%%%%%%%%%%%%%%%%%%%%%%%%%%%%%%%%%%%%%%%%%%%%%%%
\thispagestyle{empty}


\noindent\includegraphics[width=\textwidth]{cabecera}\vspace{1cm}

\noindent D. nombre tutor, profesor del departamento de nombre departamento, área de nombre área.

\noindent Expone:

\noindent Que el alumno D. \nombre, con DNI 71303809Z, ha realizado el Trabajo final de Grado en Ingeniería Informática titulado título de TFG. 

\noindent Y que dicho trabajo ha sido realizado por el alumno bajo la dirección del que suscribe, en virtud de lo cual se autoriza su presentación y defensa.

\begin{center} %\large
En Burgos, {\large \today}
\end{center}

\vfill\vfill\vfill

% Author and supervisor
\begin{minipage}{0.45\textwidth}
\begin{flushleft} %\large
Vº. Bº. del Tutor:\\[2cm]
D. nombre tutor
\end{flushleft}
\end{minipage}
\hfill
\begin{minipage}{0.45\textwidth}
\begin{flushleft} %\large
Vº. Bº. del co-tutor:\\[2cm]
D. nombre co-tutor
\end{flushleft}
\end{minipage}
\hfill

\vfill

% para casos con solo un tutor comentar lo anterior
% y descomentar lo siguiente
%Vº. Bº. del Tutor:\\[2cm]
%D. nombre tutor


\newpage\null\thispagestyle{empty}\newpage




\frontmatter

% Abstract en castellano
\renewcommand*\abstractname{Resumen}
\begin{abstract}
En este primer apartado se hace una \textbf{breve} presentación del tema que se aborda en el proyecto.
\end{abstract}

\renewcommand*\abstractname{Descriptores}
\begin{abstract}
Palabras separadas por comas que identifiquen el contenido del proyecto Ej: servidor web, buscador de vuelos, android \ldots
\end{abstract}

\clearpage

% Abstract en inglés
\renewcommand*\abstractname{Abstract}
\begin{abstract}
A \textbf{brief} presentation of the topic addressed in the project.
\end{abstract}

\renewcommand*\abstractname{Keywords}
\begin{abstract}
keywords separated by commas.
\end{abstract}

\clearpage

% Indices
\tableofcontents

\clearpage

\listoffigures

\clearpage

\listoftables
\clearpage

\mainmatter
\capitulo{1}{Introducción}

Descripción del contenido del trabajo y de la estructura de la memoria y del resto de materiales entregados.

\capitulo{2}{Objetivos del proyecto}

Este apartado explica de forma precisa y concisa cuales son los objetivos que se persiguen con la realización del proyecto. Se puede distinguir entre los objetivos marcados por los requisitos del software a construir y los objetivos de carácter técnico que plantea a la hora de llevar a la práctica el proyecto.

\capitulo{3}{Conceptos teóricos}

Este proyecto aplica la Ingeniería Informática en el ámbito de la Biomedicina, colaborando, más concretamente, con el diagnóstico por medio de imágenes en Neurología. Esto implica un proceso de investigación y aprendizaje necesarios para el entendimiento del entorno, lo que supone mayor comprensión del vocabulario y objetivos específicos, así como para el desarrollo del proyecto en sí.

En este caso, nos enfocamos en conceptos básicos bastante concretos, dentro de dos áreas de conocimiento:
\begin{itemize}
	\item Neurología
	\item Técnicas de Imagen Biomédica
\end{itemize}

En el primer caso, los conceptos neurológicos que nos interesan son los referentes al accidente (o ataque) cerebro-vascular (ACV), más frecuentemente denominado ictus, y al tejido isquémico que resulta del mismo.

En cuanto a las Técnicas de Imagen Biomédica, nos centramos en las Tomografías (Axiales) Computarizadas, conocidas comúnmente como TC o TAC. Principalmente, debemos tener claro el concepto de TC craneal, ya que es el tipo de imagen para el que este proyecto ha sido diseñado. Sin embargo, no está de más conocer algo sobre la Resonancia Magnética (MRI por sus siglas en inglés \textit{Magnetic Resonance Imaging}), ya que su comparativa con la TC da más sentido e importancia a este proyecto.

Por otro lado, en el ámbito de la Ingeniería Informática, se manejan conocimientos relacionados de una u otra manera con el Aprendizaje Automático y su desarrollo. Esta rama, quizás más conocida por su homónimo en inglés \textit{Machine Learning}, engloba conceptos teóricos de áreas como:

\begin{itemize}
	\item Inteligencia Artificial
	\item Redes Neuronales
\end{itemize}

Posteriormente, se desarrollan todos estos conceptos con más detalle.


\section{Accidente Cerebrovascular o Ictus}

El accidente (o ataque) cerebrovascular (ACV) es más comúnmente conocido como ictus, pero también mediante otras expresiones como infarto o derrame cerebral, embolia, trombosis o apoplejía \cite{wiki:acv}. No se trata de un uso preciso de dichos términos, ya que los últimos mencionados no son sinónimos, solo que pueden (o no) estar frecuentemente relacionados con la causa o tipo de ictus. Por ejemplo, el hecho de referirse a esta enfermedad cerebrovascular como ``infarto'' frente a ``derrame'' o ``hemorragia'' cerebral, se explica por los dos tipos en que se presenta:
\begin{description}
	\item [Ictus isquémico:] causado por una disminución importante en el flujo sanguíneo de una parte del cerebro, ante la oclusión de una arteria, provocando que el tejido cerebral de la zona no reciba la sangre y oxígeno necesarios para su normal funcionamiento. Se calcula que alrededor del 87\% de los ictus son isquémicos \cite{wiki:stroke}.	
	\imagen{ictus_isquemico}{Esquema representativo de un ictus isquémico.}
	
	\item [Ictus hemorrágico:] debido a la rotura de un vaso cerebral, resultando en un sangrado que aumenta la presión sobre el cerebro en esa zona, afectando a su correcto funcionamiento. En algún caso, se pueden combinar, sin saber con certeza cuántos ACV de tipo hemorrágico tienen origen isquémico, ya que se puede producir un sangrado en zonas de isquemia \cite{wiki:stroke}.
	\imagen{ictus_hemorragico}{Esquema representativo de un ictus hemorrágico.}
\end{description}

En ambos casos, bien por la falta de riego sanguíneo, bien por la presión que ejerce la hemorragia, se produce un daño en el tejido cerebral, pudiendo tener resultados fatales. Es decir, cualquiera que sea el tipo de ictus, como consecuencia se produce una lesión en las células cerebrales de un área del cerebro, las cuales pueden llegar a morir. Esto implica una disminución de la funcionalidad cerebral en esa zona, que puede suponer una discapacidad total o parcial del paciente e incluso su muerte en el peor de los casos.

¿Cuál es el papel de las técnicas de neuroimagen en la identificación y diagnóstico de un ataque cerebrovascular? Principalmente, se trata de evaluar y concretar lo siguiente, en este orden:

\begin{enumerate}
	\item Confirmar la existencia de un ACV.
	\item Determinar el tipo de ictus.
	\item Delimitar la región del cerebro afectada.
\end{enumerate} 


\section{Tomografía Computarizada (TC/TAC)}




\subsection{TC o TAC craneal}



\imagen{ictus_TC}{Imágenes de TC craneal mostrando lesiones de ictus isquémico (A) y hemorrágico (B). Se ve el tejido isquémico oscuro y el sangrado blanco.}


\subsection{Alternativa a la Resonancia Magnética (MRI)}

¡¡¡EXPLICACIÓN DE MRI!!!

Es importante destacar que este tipo de tecnología para el diagnóstico mediante imagen se encuentra menos disponible que la de CT, siendo sus costes considerablemente más elevados, tanto en términos económicos, como de tiempo \cite{MRIvsTC}. En la práctica, esto resulta clave, puesto que el ictus es una enfermedad cerebrovascular en la que el tiempo de respuesta por parte de los expertos puede ser clave, no solo a la hora de determinar el grado de discapacidad que supondrá para el paciente, sino, en algunos casos, puede incluso marcar la diferencia entre la vida y la muerte. 

Además, no todos los centros médicos pueden permitirse contar con la tecnología y personal experto de MRI, mientras que la de TC es más asequible. Junto a esto, hay que tener en cuenta el hecho de que en algunas personas, por motivos de claustrofobia o de llevar cierto tipo de implante, no puede llevarse a cabo este tipo de prueba magnética  \cite{MRIvsTC}.

Lo que también es cierto, es que resulta más sencillo para el personal médico experto evaluar y reconocer con mayor precisión y seguridad el tejido isquémico, mediante una resonancia magnética \cite{MRIvsTC}. Sin embargo, esto sólo refuerza el hecho de que el desarrollo de una tecnología como esta, por la que apuesta este proyecto, podría suponer un antes y un después, llegando más rápidamente a gran cantidad de pacientes, a través de un mayor número de centros médicos.


\section{Inteligencia Artificial}




\section{Minería de Datos}




\section{Redes Neuronales}









\indexspace
\indexspace
\indexspace
\indexspace
Algunos conceptos teóricos de \LaTeX \footnote{Créditos a los proyectos de Álvaro López Cantero: Configurador de Presupuestos y Roberto Izquierdo Amo: PLQuiz}.


\subsubsection{Subsubsecciones}

Y sub - sub - secciones. 


\section{Listas de items}

Existen tres posibilidades:

\begin{itemize}
	\item primer item.
	\item segundo item.
\end{itemize}

\begin{enumerate}
	\item primer item.
	\item segundo item.
\end{enumerate}

\begin{description}
	\item[Primer item] más información sobre el primer item.
	\item[Segundo item] más información sobre el segundo item.
\end{description}
	
\begin{itemize}
\item 
\end{itemize}

\section{Tablas}

Igualmente se pueden usar los comandos específicos de \LaTeX o bien usar alguno de los comandos de la plantilla.

\tablaSmall{Herramientas y tecnologías utilizadas en cada parte del proyecto}{l c c c c}{herramientasportipodeuso}
{ \multicolumn{1}{l}{Herramientas} & App AngularJS & API REST & BD & Memoria \\}{ 
HTML5 & X & & &\\
CSS3 & X & & &\\
BOOTSTRAP & X & & &\\
JavaScript & X & & &\\
AngularJS & X & & &\\
Bower & X & & &\\
PHP & & X & &\\
Karma + Jasmine & X & & &\\
Slim framework & & X & &\\
Idiorm & & X & &\\
Composer & & X & &\\
JSON & X & X & &\\
PhpStorm & X & X & &\\
MySQL & & & X &\\
PhpMyAdmin & & & X &\\
Git + BitBucket & X & X & X & X\\
Mik\TeX{} & & & & X\\
\TeX{}Maker & & & & X\\
Astah & & & & X\\
Balsamiq Mockups & X & & &\\
VersionOne & X & X & X & X\\
} 


\capitulo{4}{Técnicas y herramientas}

En este apartado se pretende presentar las herramientas de desarrollo utilizadas para llevar a cabo el proyecto, junto con las metodologías correspondientes.


\section{Python}

Python es un lenguaje de programación interpretado, con propósito general, tipado dinámicamente y con independencia de plataformas. Soporta programación orientada a objetos, así como imperativa y en menor medida también se utiliza para el desarrollo web o en programación funcional. \cite{python}

Su relativamente reciente expansión y popularidad se explican sobretodo por su simplicidad, ya que su sintaxis favorece un código legible, junto con su versatilidad, rapidez de desarrollo y la gran cantidad de librerías que contiene. Se trata además de un software gratuito, incluso para propósitos empresariales.

En concreto, en este trabajo se ha utilizado Python 3. Todo lo necesario sobre esta herramienta figura en <https://www.python.org/>.


\section{LaTeX}

Latex es un sistema de composición de textos de código abierto, formado por un amplio conjunto de macros de TeX, para facilitar la creación de documentos con una alta calidad tipográfica. Esto hace que su uso sea muy popular y conveniente para este tipo de documentación, así como tesis, artículos académicos o libros científicos. \cite{wiki:latex}

Puede encontrarse todo lo necesario sobre LaTex en <https://www.latex-project.org/>.


\section{Git}

Git es un sistema de control de versiones distribuido para código abierto, cuyo cometido principal es rastrear los cambios realizados durante el desarrollo de software, a la par que coordinar el trabajo entre programadores. Esto permite volver a un punto anterior del desarrollo, si fuera necesario por algún error, o recurrir a cualquier parte de la información recopilada en caso de pérdidas, así como llevar un registro continuado de cuándo y por quién se realizan modificaciones. \cite{wiki:git}

Actualmente, Git es uno de los sistemas más utilizados con este propósito y proporciona todas las herramientas necesarias para llevar el
control de versiones de una forma simple y eficiente.

Esta herramienta se encuentra a través de la URL <https://git-scm.com/>.


\section{Github}

Github es una plataforma subsidiaria de Microsoft de código abierto, para su desarrollo colaborativo. La utilización de este servicio de alojamiento basado en web permite contar con un repositorio de este proyecto, así como con todas las funciones de control de versiones distribuidas y administración de código fuente de Git, además de sus propias funciones agregadas. \cite{wiki:github}

Además, permite la integración con otras herramientas muy útiles para la planificación de tareas del proyecto, como Zenhub, la cual se detallará a continuación.

URL de la herramienta: <https://github.com/>.


\section{Zenhub}

Zenhub es una herramienta de gestión de proyectos integrada con GitHub nativamente, lo que supone que no son necesarias más configuraciones a la hora de incorporarla al repositorio del proyecto.

Las aportaciones que brinda ZenHub a la hora de planificar las distintas fases y tareas del proyecto es de gran ayuda, ofreciendo, por ejemplo, añadir de una manera muy sencilla e intuitiva un coste a las distintas \textit{issues} generadas, así como un tablero Kanban integrado y una serie de diagramas (Burndown y Velocity tracking) que se corresponden con las historias de usuarios y tareas definidas. Así, permite determinar puntos de historia y, en general, planificar de un modo más eficaz nuestros \textit{sprints} o fases temporales constituidas por cierta agrupación de tareas asignadas.

Al tratarse de una herramienta integrada en GitHub, se puede encontrar dentro de su misma página web, más concretamente en la siguiente URL: <https://github.com/marketplace/zenhub>.


\section{Jupyter}

El Proyecto Jupyter se crea para desarrollar \textit{software} de código y estándares abiertos, así como servicios para la computación interactiva en gran cantidad de lenguajes de programación. \cite{wiki:jupyter}


\subsection{Jupyter Notebook}

Se trata de un entorno computacional interactivo basado en la web, para la creación de documentos JSON referidos como \textit{Jupyter notebooks} y de extensión ".ipynb". Este tipo de ficheros siguen un esquema versionado y contienen una serie de celdas ordenadas de entrada/salida, las cuales cuentan principalmente con código y texto (usando \textit{Markdown}), aunque también pueden incluir matemáticas, gráficos, etc. \cite{wiki:jupyter}

Cuenta con una visualización simplificada y fácilmente accesible en la web, gracias a la librería \textit{nbconvert}. Ésta viene proporcionada a través de NbViewer y lleva una URL a cualquier \textit{notebook} público disponible, para convertirlo sobre la marcha en HTML y, así, mostrarlo al usuario.


\subsection{Colaboratory}

Colaboratory, también conocido como "Colab", es un entorno libre de \textit{Jupyter Notebook}, cuya ejecución se lleva a cabo en la nube y que guarda sus \textit{notebooks}
o ficheros de código en Google Drive. Desde septiembre de 2018, Colaboratory solo soporta los núcleos de Python 2 y Python 3, ningún otro núcleo de Jupyter. \cite{wiki:jupyter}

 Se trata de una opción interesante, sobretodo en este caso a la hora de hacer pruebas y comparar distintos \textit{softwares}, para evitar tener que hacer todas esas instalaciones en nuestra máquina, tanto a nivel de espacio como de conflictos de librerías, versiones, etc. Por otra parte, en este proyecto se hacía necesario el uso de GPU, frente a una CPU que se podía quedar corta de potencia, lo cual es precisamente una de las ventajas que ofrece Google Colab en la nube. A través de la siguiente URL se accede a estos \textit{notebooks}: <https://colab.research.google.com/>.


\section{SSH}

Se conoce como \textit{Secure SHell} (SSH) tanto al protocolo, como al programa que lo implementa, con la función principal del acceso remoto a un servidor, por medio de un canal seguro con la información cifrada. \cite{wiki:ssh}

Como se mencionaba previamente, para llevar a cabo este proyecto se hace necesario un procesamiento por parte de GPU. Para ello, disponemos de una máquina a la que acceder a través de un servidor, en este modo remoto. En concreto, se accede a través de una línea de comandos con la siguiente estructura: ssh -X -p 22 usuario@IPservidor. Así, se indica que se van a ejecutar programas con interfaz gráfica  (parámetro "-X"), utilizando el puerto 22, según el segundo parámetro.


\section{FTP}

Siguiendo con la utilización de esa máquina más potente, esto supone que se necesite hacer uso de un Protocolo de Transferencia de Archivos, conocido como FTP por sus siglas en inglés (\textit{File Transfer Protocol}), para el intercambio de archivos entre dicha máquina y el equipo personal desde el que se accede a ella en cualquier momento.

Se trata de un protocolo de red basado en la arquitectura cliente-servidor, cuya función es transferir archivos entre sistemas conectados a una red TCP (\textit{Transmission Control Protocol}), independientemente de sus sistemas operativos. \cite{wiki:ftp}

En este caso, para llevar a cabo estas tareas de transferencia de archivos, decidimos llevar a cabo directamente en el router la configuración de un servidor FTP y un usuario, a través del cual acceder como cliente desde cualquier terminal.



\capitulo{5}{Aspectos relevantes del desarrollo del proyecto}

Este apartado pretende recoger los aspectos más interesantes del desarrollo del proyecto, comentados por los autores del mismo.
Debe incluir desde la exposición del ciclo de vida utilizado, hasta los detalles de mayor relevancia de las fases de análisis, diseño e implementación.
Se busca que no sea una mera operación de copiar y pegar diagramas y extractos del código fuente, sino que realmente se justifiquen los caminos de solución que se han tomado, especialmente aquellos que no sean triviales.
Puede ser el lugar más adecuado para documentar los aspectos más interesantes del diseño y de la implementación, con un mayor hincapié en aspectos tales como el tipo de arquitectura elegido, los índices de las tablas de la base de datos, normalización y desnormalización, distribución en ficheros3, reglas de negocio dentro de las bases de datos (EDVHV GH GDWRV DFWLYDV), aspectos de desarrollo relacionados con el WWW...
Este apartado, debe convertirse en el resumen de la experiencia práctica del proyecto, y por sí mismo justifica que la memoria se convierta en un documento útil, fuente de referencia para los autores, los tutores y futuros alumnos.

\capitulo{6}{Trabajos relacionados}

Este apartado sería parecido a un estado del arte de una tesis o tesina. En un trabajo final grado no parece obligada su presencia, aunque se puede dejar a juicio del tutor el incluir un pequeño resumen comentado de los trabajos y proyectos ya realizados en el campo del proyecto en curso. 

\capitulo{7}{Conclusiones y Líneas de trabajo futuras}

Todo proyecto debe incluir las conclusiones que se derivan de su desarrollo. Éstas pueden ser de diferente índole, dependiendo de la tipología del proyecto, pero normalmente van a estar presentes un conjunto de conclusiones relacionadas con los resultados del proyecto y un conjunto de conclusiones técnicas. 
Además, resulta muy útil realizar un informe crítico indicando cómo se puede mejorar el proyecto, o cómo se puede continuar trabajando en la línea del proyecto realizado. 



\bibliographystyle{plain}
\bibliography{bibliografia}

\end{document}
